\section{Arytmetyka Peana i Robinsona}\label{sec:arytmetyka-peana-robinsona}
Zanim podamy aksjomaty, wyjaśnijmy, po co one są. Chcemy zapisać w języku $\mathcal L_A$ (Def.~\ref{res:3:1}) na tyle dużo prawd o liczbach naturalnych, by teoria potrafiła „liczyć” --- a jednocześnie zrobić to skończoną, mechanicznie sprawdzalną listą reguł. Arytmetyka Robinsona $Q$ to taki możliwie minimalny zestaw: opisuje tylko podstawowe własności zera, następnika, dodawania, mnożenia i porządku.

\begin{definicja}[Arytmetyka Robinsona $Q$]\label{res:4:1}
\emph{Arytmetyką Robinsona} $Q$ nazywamy teorię o następujących ośmiu aksjomatach. Nazwy (Q1)--(Q8) są stałe, niezależne od numeracji wzorów:
\begin{align*}
\forall x\ S(x)\ne0, &\tag{Q1}\\
\forall x\,\forall y\,(S(x)=S(y)\to x=y), &\tag{Q2}\\
\forall x\,(x\ne0\to\exists y\ x=S(y)), &\tag{Q3}\\
\forall x\ x+0=x, &\tag{Q4}\\
\forall x\,\forall y\ x+S(y)=S(x+y), &\tag{Q5}\\
\forall x\ x\cdot0=0, &\tag{Q6}\\
\forall x\,\forall y\ x\cdot S(y)=x\cdot y+x, &\tag{Q7}\\
\forall x\,\forall y\ (x<y\leftrightarrow\exists z\,(z+S(x)=y)). &\tag{Q8}
\end{align*}
Przyjmujemy $x\le y$ jako skrót $x<y\lor x=y$. Kolejność składników w (Q8) ma znaczenie: $Q$ nie dowodzi przemienności dodawania. Ta definicja zapewnia, że odcinek poniżej standardowego numerału jest dokładnie skończonym zbiorem mniejszych numerałów.
\end{definicja}

\begin{uwaga}[Znaczenie każdego aksjomatu --- po ludzku]\label{res:4:2}
Czytajmy $S(x)$ jako „$x+1$”, a $0$ jako zero. Wtedy:
\begin{itemize}
\item \textup{(Q1)} $S(x)\ne0$: \emph{zero nie jest niczyim następnikiem} --- żadna liczba $+1$ nie daje $0$. To odróżnia $\mathbb N$ od np. liczb modulo, gdzie „zawinięcie” pozwoliłoby wrócić do zera.
\item \textup{(Q2)} $S(x)=S(y)\to x=y$: \emph{następnik jest różnowartościowy} --- jeśli $x+1=y+1$, to $x=y$. Gwarantuje, że kolejne liczby się nie „sklejają”.
\item \textup{(Q3)} $x\ne0\to\exists y\ x=S(y)$: \emph{każda niezerowa liczba ma poprzednik} --- nie ma „dziur” ani dziwnych elementów bez $-1$.
\item \textup{(Q4)} $x+0=x$: dodanie zera nic nie zmienia (baza rekurencji dodawania).
\item \textup{(Q5)} $x+S(y)=S(x+y)$: \emph{definiuje dodawanie przez rekurencję} --- dodać $y+1$ to dodać $y$ i potem $+1$. Aksjomaty (Q4)--(Q5) wspólnie mówią, jak liczyć $x+y$ krok po kroku.
\item \textup{(Q6)} $x\cdot0=0$: mnożenie przez zero daje zero (baza).
\item \textup{(Q7)} $x\cdot S(y)=x\cdot y+x$: \emph{definiuje mnożenie przez rekurencję} --- pomnożyć przez $y+1$ to pomnożyć przez $y$ i dodać $x$.
\item \textup{(Q8)} $x<y\leftrightarrow\exists z\,(z+S(x)=y)$: w modelu standardowym $y$ otrzymujemy, dodając $x+1$ do pewnego $z\in\mathbb N$, dokładnie w kolejności $z+S(x)$. Nie zakłada ogólnych praw porządku, których $Q$ mogłaby nie dowodzić.
\end{itemize}
\end{uwaga}

\begin{przyklad}[$Q$ liczy konkretne przykłady]\label{res:4:3}
Z aksjomatów (Q4)--(Q5) $Q$ dowodzi np. $\overline2+\overline1=\overline3$: istotnie $\overline2+S(0)=S(\overline2+0)=S(\overline2)=\overline3$. Ogólnie $Q$ poprawnie dowodzi każdej prawdziwej równości i nierówności między \emph{konkretnymi} numerałami. To wystarczy do reprezentowalności (Tw.~\ref{res:7:6}).
\end{przyklad}

\begin{lemat}[Skończone odcinki i rachunek na numerałach]\label{res:4:4}
Dla każdego ustalonego standardowego $n$ teoria $Q$ dowodzi
\begin{align}
\forall x\,\Bigl(x<\overline n\leftrightarrow\bigvee_{j<n}(x=\overline j)\Bigr),\label{eq:4:1}\\
\forall x\,(x<\overline n\lor x=\overline n\lor\overline n<x).\label{eq:4:2}
\end{align}
Dla $n=0$ alternatywa po prawej w~\eqref{eq:4:1} jest pusta i oznacza fałsz. Ponadto $Q$ rozstrzyga równości i nierówności między termami domkniętymi.
\end{lemat}

\begin{proof}
Z (Q4)--(Q7) obliczamy każdy term domknięty, rozwijając dodawanie i mnożenie skończenie wiele razy. Różność dwóch różnych numerałów wynika z (Q1)--(Q2).

Z (Q3) i prawa wyłączonego środka otrzymujemy $x=0\lor\exists u\ x=S(u)$. W drugiej gałęzi stosujemy ten sam podział do $u$: albo $u=0$ i $x=\overline1$, albo $u=S(v)$ i $x=S^2(v)$. Powtarzamy to dokładnie $n+1$ razy, otrzymując w $Q$
\[
x=\overline0\lor x=\overline1\lor\cdots\lor x=\overline n\lor\exists u\ x=S^{n+1}(u).
\]
Zapis $S^k(u)$ oznacza $k$ kolejnych zastosowań następnika. Jeśli $x=S^{n+1}(u)$, to $z+S(x)=S^{n+2}(z+u)\ne\overline n$ na mocy (Q1)--(Q2) i (Q5): domniemana równość po $n$ skróceniach następnika dawałaby $S^2(z+u)=0$. Zatem $x\not<\overline n$. Dla $x=\overline j$ z $j\le n$ zachodzi $z+S(\overline j)=S^{j+1}(z)$. Przy $j<n$ świadkiem równości z $\overline n$ jest $z=\overline{n-j-1}$. Przy $j=n$ taka równość po $n$ skróceniach dawałaby $S(z)=0$, więc jest niemożliwa. To dowodzi obu kierunków~\eqref{eq:4:1}. W ostatnim członie alternatywy $u+S(\overline n)=S^{n+1}(u)=x$, więc $\overline n<x$; pozostałe gałęzie dają~\eqref{eq:4:2}. Jest to indukcja w metateorii budująca osobny skończony dowód dla każdego $n$, a nie zastosowanie nieistniejącego w $Q$ schematu indukcji.
\end{proof}

\begin{uwaga}[Jak połączyć skończenie wiele sprawdzeń]\label{res:4:5}
Niech $n$ będzie ustaloną liczbą standardową, a $\theta(x)$ dowolną formułą. Jeśli $Q$ dowodzi wszystkich $\theta(\overline0),\ldots,\theta(\overline{n-1})$, to dowodzi też $\forall x<\overline n\,\theta(x)$. Ten ostatni zapis jest skrótem $\forall x(x<\overline n\to\theta(x))$. Istotnie, z $x=\overline j$ oraz $\theta(\overline j)$ otrzymujemy $\theta(x)$ przez podstawienie równych termów. Łączymy skończenie wiele takich implikacji przez alternatywę i stosujemy~\eqref{eq:4:1}. Następnie generalizujemy $x$. Dla $n=0$ teza wynika z $\forall x\,\neg(x<0)$. Tak powstaje \emph{jeden skończony dowód dla ustalonej granicy} $n$, a nie dowód $\forall x\,\theta(x)$. Tego kroku użyjemy między innymi w dowodzie Rossera.
\end{uwaga}

\begin{definicja}[Formuły ograniczone i $\Sigma_1$]\label{res:4:6}
Formuła $\Delta_0$ powstaje z formuł atomowych przez spójniki oraz kwantyfikatory $\exists x<t$, $\forall x<t$, gdzie $t$ nie zawiera $x$. Są to skróty $\exists x(x<t\land\cdots)$ i $\forall x(x<t\to\cdots)$. Kwantyfikatory z $x\le t$ również są ograniczone: oznaczają odpowiednio $(\exists x<t\,\varphi(x))\lor\varphi(t)$ oraz $(\forall x<t\,\varphi(x))\land\varphi(t)$, po podstawieniu bez przechwytywania zmiennych. Nie potrzeba tu żadnych ogólnych praw porządku. Formuła $\Sigma_1$ ma postać $\exists u_1\cdots\exists u_r\,\delta$, gdzie $\delta$ jest $\Delta_0$; dopuszczamy $r=0$.
\end{definicja}

\begin{lemat}[$\Sigma_1$-zupełność $Q$]\label{res:4:7}
$Q$ dowodzi każdego prawdziwego zdania $\Sigma_1$ oraz negacji każdego fałszywego zdania $\Delta_0$.
\end{lemat}
\begin{proof}
Po podstawieniu numerałów termy obliczamy według Lematu~\ref{res:4:4}. Każdy kwantyfikator ograniczony przez numerał zastępujemy skończoną alternatywą lub koniunkcją za pomocą~\eqref{eq:4:1}. Indukcja po budowie formuły daje dla każdej zamkniętej $\Delta_0$ dowód jej prawdziwej strony: samej formuły albo jej negacji. Jeżeli $\exists\overline u\,\delta(\overline u)$ jest prawdziwa, wybieramy rzeczywistych świadków $\overline m$, dowodzimy $\delta(\overline m)$ w $Q$ i wprowadzamy kwantyfikatory egzystencjalne. Nie otrzymujemy stąd rozstrzygania fałszywych zdań $\Sigma_1$.
\end{proof}

\begin{uwaga}[Czego $Q$ nie potrafi]\label{res:4:8}
$Q$ jest celowo słaba: \emph{nie} dowodzi ogólnych praw z kwantyfikatorem, np. przemienności $\forall x\,\forall y\ x+y=y+x$ ani $\forall x\ 0+x=x$. Powód: brakuje jej \emph{indukcji}. Mimo to jej słabość jest zaletą --- im słabsza teoria, dla której udowodnimy niezupełność, tym mocniejszy wynik (bo obejmie automatycznie wszystkie silniejsze).
\end{uwaga}

\begin{definicja}[Arytmetyka Peana PA]\label{res:4:9}
\emph{Arytmetyką Peana} $\mathrm{PA}$ nazywamy teorię powstałą z $Q$ (aksjomaty (Q1)--(Q7), z (Q8) jako definicją $<$) przez dodanie \emph{schematu indukcji}: dla każdej formuły $\varphi(x,\overline y)$ języka $\mathcal L_A$ przyjmujemy aksjomat
\[
\forall\overline y\,\bigl[(\varphi(0,\overline y)\land\forall x\,(\varphi(x,\overline y)\to\varphi(S(x),\overline y)))\to\forall x\,\varphi(x,\overline y)\bigr].
\]
\end{definicja}

\begin{uwaga}[Co mówi schemat indukcji]\label{res:4:10}
To formalna wersja „zasady domina”: jeśli własność $\varphi$ zachodzi dla $0$ (\emph{baza}) i z jej zachodzenia dla $x$ wynika zachodzenie dla $x+1$ (\emph{krok}), to zachodzi dla \emph{wszystkich} liczb. Słowo „schemat” oznacza, że jest to \emph{nieskończenie wiele} aksjomatów --- po jednym dla każdej formuły $\varphi$. Mimo nieskończoności zbiór ten jest \emph{rozstrzygalny} (można mechanicznie sprawdzić, czy dane zdanie ma kształt aksjomatu indukcji), więc PA pozostaje efektywnie aksjomatyzowalna --- co jest niezbędne dla twierdzeń Gödla.
\end{uwaga}

\begin{uwaga}\label{res:4:11}
$Q$ jest bardzo słaba (np. nie dowodzi przemienności dodawania), lecz wystarcza do reprezentowania wszystkich funkcji obliczalnych w sensie obliczeń na ustalonych numerałach (Tw.~\ref{res:7:6}). Pierwsze twierdzenie, w wersji Rossera, obejmuje każde niesprzeczne efektywnie aksjomatyzowalne rozszerzenie $Q$. Drugie twierdzenie dowodzimy tutaj dla rozszerzeń PA; sama reprezentowalność nie zastępuje warunków wyprowadzalności.
\end{uwaga}

% Rozdziały 5--6, strony 9--13.
